% Cambia la opción: memoria | deber | apuntes
\documentclass[memoria]{uc3mmaster}

\titulo{Título de la memoria}
\subtitulo{Subtítulo opcional}
\asignatura{Nombre de la asignatura}
\tipo{Memoria}                % p. ej. Práctica 3, Apuntes, Trabajo final
\curso{Curso 2026--2027}
\fecha{Octubre de 2026}
% \logo{logo-uc3m.png}        % opcional: pon aquí tu logo (no incluido)
\autor{Nombre Apellidos}{NIA}
\autor{Nombre Apellidos}{NIA}

\begin{document}
\makeportada

\section{Introducción}
Texto con fórmulas: $h^{*}(s)=\min_{\pi}\sum_{o\in\pi}c(o)$.

\begin{definicion}[Heurística admisible]
$h(s)\le h^{*}(s)$ para todo estado $s$.
\end{definicion}
\begin{teorema}[A$^*$]  A$^*$ con $h$ admisible es óptimo. \end{teorema}
\begin{ejemplo} Un ejemplo breve. \end{ejemplo}
\begin{nota} Aviso o comentario. \end{nota}
\begin{resultado}[Resultado]  $h=6$. \end{resultado}

\begin{pseudocodigo}{Búsqueda}
\pl $\mathit{Abiertos}\gets\{s_0\}$
\pl \kw{mientras} $\mathit{Abiertos}\neq\emptyset$ \kw{hacer}
\pl \ind extraer $n$ con menor $f$
\end{pseudocodigo}

\begin{table}[h]\centering
\begin{tabular}{@{}lcc@{}}\toprule A & B & C\\\midrule 1 & 2 & 3\\\bottomrule\end{tabular}
\caption{Ejemplo de tabla.}
\end{table}

\begin{thebibliography}{9}\footnotesize
\bibitem{x} Autor. Título. Editorial, 2026.
\end{thebibliography}

\makecontraportada   % quítala si no quieres contraportada
\end{document}
